\chapter{Partizip II and the Perfect Tenses}

The Perfekt combines a present-tense form of \textit{haben} or \textit{sein}
with the Partizip II. The auxiliary is finite and normally occupies position
two; the participle closes the sentence bracket.

\begin{center}
\textit{Ich \textbf{habe} den Brief gestern \textbf{geschrieben}.}
\qquad
\textit{Wir \textbf{sind} früh \textbf{abgefahren}.}
\end{center}

\section{Forming the Partizip II}

The following patterns cover most verbs. Strong and mixed forms must still be
learned as principal parts because their vowel or stem may change.

\begin{center}
\small
\rowcolors{2}{referenceBlueBg}{white}
\begin{tabularx}{\textwidth}{@{}>{\bfseries}l l l X@{}}
\toprule
Verb type & Pattern & Infinitive & Partizip II \\
\midrule
Weak & \textit{ge-} + stem + \textit{-t} & machen & gemacht \\
Weak in \textit{-d/-t} & \textit{ge-} + stem + \textit{-et} & arbeiten & gearbeitet \\
Strong & often \textit{ge-} + changed stem + \textit{-en} & schreiben & geschrieben \\
Mixed & stem change + weak ending & bringen & gebracht \\
Separable & prefix + \textit{ge-} + verb & anrufen & angerufen \\
Inseparable prefix & no \textit{ge-} & verstehen & verstanden \\
Verb in \textit{-ieren} & no \textit{ge-}, ending \textit{-t} & studieren & studiert \\
\bottomrule
\end{tabularx}
\end{center}

The common inseparable prefixes are \textit{be-, emp-, ent-, er-, ge-, miss-,
ver-}, and \textit{zer-}: \textit{bezahlt, empfohlen, entdeckt, erzählt,
gefallen, missverstanden, verkauft, zerstört}.

\section{Choosing \textit{haben} or \textit{sein}}

Most verbs form the Perfekt with \textit{haben}. This includes transitive
verbs with an accusative object, reflexive verbs, and verbs describing an
activity or state without a change of location or condition.

\begin{center}
\small
\rowcolors{2}{referenceGreenBg}{white}
\begin{tabularx}{\textwidth}{@{}>{\bfseries}l X X@{}}
\toprule
Auxiliary & Typical use & Examples \\
\midrule
haben & most verbs; transitive, reflexive, and modal verbs
      & \textit{Ich habe gelesen. Sie hat sich beeilt.} \\
sein & change of location or state; also \textit{sein, werden}, and
       \textit{bleiben}
     & \textit{Er ist angekommen. Es ist eingeschlafen.} \\
\bottomrule
\end{tabularx}
\end{center}

Meaning can affect the auxiliary. Compare
\textit{Ich \textbf{bin} nach Berlin gefahren} (movement to a destination)
with \textit{Ich \textbf{habe} den Wagen gefahren} (transitive: driving the
vehicle). Where accepted usage varies, the individual verb table shows both
auxiliaries.

\section{Word order in the Perfekt}

\begin{center}
\small
\rowcolors{2}{referenceYellowBg}{white}
\begin{tabularx}{\textwidth}{@{}>{\bfseries}l X@{}}
\toprule
Clause type & Example \\
\midrule
Main clause & Ich \textbf{habe} das Buch \textbf{gelesen}. \\
Question & \textbf{Hast} du das Buch \textbf{gelesen}? \\
Subordinate clause & \ldots{}, weil ich das Buch \textbf{gelesen habe}. \\
With negation & Ich \textbf{habe} das Buch nicht \textbf{gelesen}. \\
\bottomrule
\end{tabularx}
\end{center}

\section{The Plusquamperfekt}

The Plusquamperfekt places an event before another past event. It uses the
Präteritum of the auxiliary, \textit{hatte} or \textit{war}, plus the same
Partizip II used in the Perfekt.

\begin{center}
\textit{Nachdem ich \textbf{gegessen hatte}, ging ich spazieren.}

\textit{Sie \textbf{war schon abgefahren}, als ich anrief.}
\end{center}

\section{Modal verbs in the Perfekt}

When a modal verb accompanies another infinitive, German normally uses two
infinitives instead of the modal participle. This is called the
\textit{Ersatzinfinitiv}.

\begin{center}
\textit{Ich habe lange \textbf{arbeiten müssen}.}
\qquad
not \quad \textit{Ich habe lange arbeiten gemusst.}
\end{center}

When the modal stands alone, its normal participle is possible:
\textit{Ich habe das nicht gewollt.} In everyday speech, the Präteritum is
often simpler and more natural with modal verbs:
\textit{Ich musste lange arbeiten.}

The Ersatzinfinitiv also creates an important subordinate-clause exception:
the finite auxiliary normally precedes the two infinitives:
\textit{\ldots{}, weil ich lange \textbf{habe arbeiten müssen}}.

\begin{tcolorbox}[
  colback=referenceRedBg,
  colframe=genderFeminine,
  title={Three forms to learn together},
  fonttitle=\bfseries
]
For an irregular verb, learn the infinitive, third-person Präteritum, and
complete Perfekt form together:

\begin{center}
\textbf{gehen -- ging -- ist gegangen}
\qquad
\textbf{sehen -- sah -- hat gesehen}
\end{center}
\end{tcolorbox}
